% sections/introduction.tex
\section{Introduction}

Training speed-runs fix a budget and ask how low a validation loss can go. Community records such as the
modded-nanoGPT speed-run \citep{jordan2024moddednanogpt}, built on nanoGPT \citep{karpathy2022nanogpt}, academic
``cramming'' studies \citep{geiping2023cramming,izsak2021academic}, budgeted-training analyses \citep{li2020budgeted}
and time-to-result benchmarks such as MLPerf Training \citep{mattson2020mlperf} and AlgoPerf
\citep{dahl2023algoperf} share one property: the quantity being optimised is the quality that a \emph{wall-clock}
budget buys, not the quality per step or per FLOP. In such a setting two things are scarce. A full trial costs the
whole budget on the target hardware, and the official score comes from the organisers with limited feedback. Most
practical decisions therefore rest on proxies, such as short screens, scores on public data and runs on other
hardware, whose relation to the official score is rarely measured.

This paper reports how one team measured that relation in Phase~1 of the AWS Trainium Frontier challenge
\citep{trainiumfrontier2026rules}. A language model is trained from scratch for 30~minutes of charged time on one
\texttt{trn2.3xlarge} instance and scored by validation bits per byte (\valbpb{}) on a pinned validation shard
computed by the organisers. Over 17~days and \val{n-scored} scored uploads, our official score fell from
\val{first-official} to \val{best-official} (\cref{fig:history}). Our best upload finished outside the top ten, so
this is not a paper about a winning recipe. It is a measurement study: every upload was preceded by a full-length
rehearsal on our own hardware, failures were logged as carefully as successes, and the repository recomputes the
numbers in this paper from the released data.

\paragraph{What is and is not new.} Replaying a candidate on deterministic hardware and correcting it with a
hold-out offset is standard validation practice, and a ridge regression with era, chip and seed effects is a
standard surrogate. Our contribution is the \emph{measurement}: how well that practice predicted a real
leaderboard, how large its errors and biases were, what limited it, and what a single price per optimizer step says
about the changes we tried. We think these quantities are useful to anyone running a time-budgeted contest or
benchmark, and they are rarely reported.

\begin{figure}[tbp]
  \centering
  \includegraphics[width=\linewidth]{score_history.pdf}
  \caption{\textbf{Official score per upload.} (a) The whole campaign in calendar time; the line is the best score
    so far and the shaded region is shown in (b). (b) From \upl{K50} on, in upload order. Circles set a new best;
    squares did not: four official tests of single changes (\upl{K65a}--\upl{K65d}) and two shuffle-salt draws
    (\upl{K73s6}, \upl{K82s7}). The change behind each new best is listed in \cref{tab:uploads}. Source:
    \texttt{results/official-scores.csv}.}
  \label{fig:history}
\end{figure}

\paragraph{Contributions.} The same four contributions, in this order, structure the paper.
\begin{itemize}
  \item \textbf{C1: How well deterministic rehearsal predicts a leaderboard} (\cref{sec:oracle}). A cold,
    full-budget rehearsal of the exact upload file, scored on the first \val{eval-tokens} public validation tokens,
    predicted the official score up to an offset of \val{off-mean-since} (SD \val{off-sd-since}, $n=\val{off-n-since}$).
    The offset is mostly text (the official shard is harder than the first 2M public tokens) but carries a
    per-instance component of a few $10^{-4}$. A calibration frozen on the chip-C calibration uploads predicted the
    seven final uploads with MAE \val{holdout-mae} and a positive bias on faster instances; we report how much of
    this depends on one normalised rehearsal, on salt selection and on the burn-in of the interval.
  \item \textbf{C2: A run simulator, with baselines} (\cref{sec:sim}). \ffsim{} decomposes a score into steps,
    computed by replaying the training script's charged-time schedule, and quality per step, from a ridge surrogate
    fitted on \val{q-nfit} fully specified chip runs out of \val{ds-n} harvested records. Held out pair by pair it
    beats a steps-only baseline and a majority-sign baseline on recipe changes; forward in time its knob features add
    little; on new mechanisms it is no better than chance. A GPU proxy that replays the chip's schedule is a
    feasibility demonstration.
  \item \textbf{C3: An accounting of what moved the score} (\cref{sec:price,sec:findings}). One price,
    $\kappa=\val{kappa}$ bpb per unit $\ln(\text{steps})$, converts every change into optimizer steps. With it we
    report the levers with their noise, a row-level shuffle pool (\val{rowpool-chip} on chip,
    \val{rowpool-official} officially, one pair each), a compilation-accounting trap that only cold rehearsals
    reveal, the size of seed and order noise, negative results, and a pricing of the gap to the top ten
    (\val{gap-ten-range} more effective compute) that we offer as a hypothesis about kernels, not a finding.
  \item \textbf{C4: Released artifacts} (\cref{sec:contrib}). The simulator, the run dataset, both recipes, a
    local app and a reviewed, guarded pipeline through which others can add runs to the simulator. No external
    contributions exist yet.
\end{itemize}

\paragraph{Scope.} This is one team's campaign on one contest and one hardware type, and most late effects rest on
one or a few same-seed pairs at a seed that was itself selected. We give $n$ and a noise scale for every effect, mark
numbers that exist only in campaign documents with a dagger (\docnum{}, listed in \cref{app:docfacts}), and list the
threats to validity in \cref{sec:limits}. Other teams appear only as rank-indexed scores from the public leaderboard
snapshot.

\paragraph{Roadmap.} \Cref{sec:related} places the work in context. \Cref{sec:setting} describes the contest, the
determinism the method relies on, the price of a step and our final recipe (\cref{sec:recipe}). \Cref{sec:oracle,sec:sim} present C1
and C2, \cref{sec:findings} C3, \cref{sec:contrib} C4 and \cref{sec:limits} the limitations. The appendices hold the
upload ledger (\cref{app:ledger}), oracle details (\cref{app:oracle}), the simulator internals, the
contribution pipeline, the final hyperparameters, derivations and a notation table.
